\section{Scope, Method, and Qualifications}

\subsection{Reader, question, and claim boundary}

This study is written for a reader who wants a single verse worked to the
bottom: someone who can follow an argument about verbal form without knowing
Hebrew, and who wants to see the evidence rather than be told the result. No
prior acquaintance with rabbinic literature or with the Septuagint is assumed;
every technical term is glossed where it first appears.

The governing question is whether Genesis 15:5 --- ``Look up to heaven and
number the stars if thou canst'' --- was spoken in daylight, on the internal
evidence of a chapter whose sun is going down at verse 12 and set at verse 17.
The claim under test was supplied by the maintainer and is not this study's; the
study's task was to test it, not to establish it.

The boundary of the claim is textual. This study reaches a conclusion about what
Genesis 15 asserts about its own sequence of hours. It reaches no conclusion
about the historicity of the episode, about whether the chapter derives from one
source or several, about the date of its composition, or about what a man
standing in Canaan would have seen. It does not adjudicate between the Jewish
and Christian readings of Genesis 15:6, does not treat the doctrine of
justification, and does not treat the theology of covenant.

\subsection{Terms fixed for this study}

\textbf{Daylight reading} means the reading on which Genesis 15:5 falls before
the chapter's own sunset and therefore in the hours of light. It is the claim
under test, not a conclusion assumed.

\textbf{Intervening-day reading} means the reading on which the star-promise
precedes a night that the text does not mention, with the covenant rite falling
on the following day. It is the principal rival.

\textbf{Vision frame} means the declaration \textit{bammahazeh} at 15:1 and the
question of its extent, not any particular theory of prophecy.

\textbf{Gloss} means a working English equivalence made by this study to show
what a Hebrew, Aramaic, Greek, or Latin word is doing in an argument. A gloss is
never a translation of Scripture and is never offered for recitation or
quotation as such. Where a published translation is used, its translator is
named at the point of use.

\textbf{Project synthesis} marks an editorial conclusion or proposal of this
study. The sign-as-demand reading in the synthesis section, the extension of Rav
Pappa's limitation from halakhic derivation to narrative chronology, the
observation that the above-the-firmament reading would incidentally dissolve the
daylight problem, and the classification of the tradition into five answers are
all project synthesis and are attributed to no source.

\subsection{Method and source hierarchy}

Each witness was read at its own locus, in its own language where the argument
turned on the wording, and never through a quotation in a later author. Where
that was not possible, the witness is listed in the ``reported only'' section of
the witness register and no claim rests on it.

The hierarchy applied was: the biblical text in its identified editions first;
then ancient versions and Targums; then ancient Jewish and Christian
commentary; then medieval and later commentary; then modern editorial apparatus.
A commentator's report of an earlier authority was treated as evidence about the
commentator, not about the authority.

Three collations were performed. The Hebrew forms at Genesis 15:5, 12, and
17--18 were verified against the identified pages of the Kittel--Beer 1909
printing. The Latin of Genesis 15 was read from a page image of the Hetzenauer
printing rather than from a transcription. Abarbanel's
seventh question, on which the central historical finding of this study depends,
was read in two independent digital transcriptions of the Warsaw 1862 printing
and found to agree word for word.

\subsection{Limitations}

\textbf{Text-critical.} The Septuagint claims in this study are claims about the
text Swete prints and the witnesses in his apparatus. Rahlfs' edition, the
G\"ottingen \textit{Genesis}, and Wevers' notes on the Greek text were not
consulted. The reading \textit{ekei} at 15:18 is reported as the reading of that
edition and is explicitly qualified by Augustine's contrary Old Latin witness.
The Kittel--Beer pages verify the forms printed at the controlling loci, but
that printing was not collated against a manuscript facsimile or a modern
critical Hebrew edition.

\textbf{Not reached.} The Genesis Apocryphon and any other Qumran
witness to Genesis 15; Targum Neofiti and the Fragment Targums; the Samaritan
Pentateuch; Jerome's \textit{Hebrew Questions on Genesis}; the Glossa ordinaria;
and the sixteenth-century Catholic commentators a Lapide names. Any of the
first three could contain a
temporal marker between 15:6 and 15:9 that would settle the question against the
claim, and the synthesis says so.

\textbf{Newly bounded witnesses.} A dated Syriac web delivery was inspected,
but its textual base is unidentified; it supports only a claim about that
delivery, not the Peshitta tradition. \textit{Jubilees} 14 was inspected as
rewritten Scripture rather than as a textual version. Ambrose, Chrysostom, and
Ephrem supply exact-locus positive reception evidence; Theodoret supplies only
a local negative. Their deliveries still await reusable normalization.

\textbf{Optical character recognition.} The texts of Cornelius a Lapide, of the
1609 Douay Old Testament, of Haydock, of Swete, of Philo in Yonge's translation,
and of Josephus in Whiston's were read in digital text layers of scanned
printings. Long-\textit{s} orthography and ligatures in the a Lapide and 1609
Douay texts were normalized when quoting; no word was altered. Passages quoted
from these witnesses were read in enough surrounding context to establish that
the sense is continuous, but were not collated against page images. This is a
real limitation and it bears most heavily on the a Lapide quotations, which are
short and decisive.

\textbf{Language.} Hebrew, Aramaic, Greek, and Latin are romanized or quoted
without vowel points and without diacritical marks throughout, for reasons of
typesetting rather than scholarship. Readers who need pointed or accented text
should consult the editions named in the references.

\textbf{Negative results.} Every negative result in this study is bounded by an
enumerated list of witnesses actually read. None of them is a claim that a
difficulty was never noticed by anyone; each is a claim that it was not noticed
in the named places.

\subsection{Rights and quotation boundary}

Scripture in English is quoted only from the registered public-domain
Douay--Rheims in Challoner's revision, and in the Vulgate numbering that witness
uses. This study composes no translation of Scripture and offers none.

Third-party translations are quoted briefly, always with the translator named,
and they retain their own rights status: Friedl\"ander on Maimonides, Dods on
Augustine, Yonge on Philo, Whiston on Josephus, Etheridge on Pseudo-Jonathan,
and Rosenbaum and Silbermann on Rashi are all in the public domain. Chavel on
Nahmanides, Munk on Radak, Chizkuni, Rabbeinu Bahya, and the Tur, and the
Sefaria Midrash Rabbah are published under Creative Commons Attribution licences
and are quoted briefly with attribution. The William Davidson Talmud and
Strickman and Silver on Ibn Ezra are published under Creative Commons
Attribution-NonCommercial licences; they were consulted only as controls on this
study's own glosses of the Aramaic and Hebrew, and their wording is deliberately
not reproduced anywhere in this study. No third-party translation is offered
under this project's licence.

Hebrew, Aramaic, Greek, and Latin source wording is quoted from editions
identified at the point of use. The Latin page image of the Hetzenauer printing
and the checked transcription made from it are retained under the rights basis
recorded with those artifacts. Project-created glosses, tables, and synthesis are
project content.

\subsection{Review state}

Completed for this revision: profile routing to the scriptural-studies section of
the discursive-articles profile; claim-driven witness selection; direct reading
of every witness listed in the register at the locus recorded there; independent
collation of the Abarbanel passage against a second transcription; a check that
no verse quoted in English falls among the recorded divergences between the
vendored Douay--Rheims and the 1899 American edition; source-library validation;
multipass build with a clean log; and raster review of every rendered page.

Outstanding: independent review by a Hebraist, a Septuagint specialist, a
rabbinics scholar, a patristics scholar, and a Catholic exegete. The
text-critical limitations above have not been closed. No imprimatur, nihil
obstat, or other ecclesiastical approval is claimed, and none of the conclusions
here has any magisterial standing. This is a study aid.
